\chapter{Inspection and recording}
\label{ch:s-inspection}

\begin{notebox}
Teach the inspection points in NETRA first (\menupath{Substation › Inspection points}), on the
stops of the route under test, and set the day's air temperature on the robot.
\end{notebox}

\begin{tests}
  I1 & Run the route; at each stop, watch the camera on the H12. & The camera turns to every
    target; NETRA gets a result for every point, with the thermal and visible images. & \result \\
  I2 & Measure one target with the reference thermometer at the same time. & The robot's reading
    within \SI{\pm 2}{\celsius} or \SI{\pm 2}{\percent} of the reference. Robot \blank, reference
    \blank & \result \\
  I3 & Put the heat source on a target of a point with three regions (phases), about \SI{10}{K} above
    the others. & The finding is at least warning (similar criterion), the hot region marked. &
    \result \\
  I4 & A partial discharge point, with a source or a known corona. & The class and the severity in
    NETRA match the source (corona: attention). & \result \\
  I5 & \ui{Inspect} on the H12 at a stop point. & The point is inspected where the robot stands; the
    result on the H12 and in NETRA. & \result \\
  I6 & During the patrol, the recorder. & \ui{REC} on the H12 and in NETRA; a session folder on the
    perception host with video and thermal frames. & \result \\
  I7 & Recorder on and off from NETRA (\menupath{Live control}). & It starts and stops; the session
    is kept. & \result \\
\end{tests}
